\scanpage{90}

\[ c_{\varpi_*}(\underset{\substack{\Vert\\\omega_{\unsure{0|1}}}}{\lambda\delta})
 =\underbrace{\delta\,d\lambda}_{\text{poids 2}}+
 \deleted{\unsure{\lambda}}\bigl[\underset{\substack{\Vert\\0}}{[\varpi_2,\lambda\delta]}+
 \underset{\substack{\Vert\\\underbrace{+\varpi_1(\lambda\delta)}_{\text{poids 2}}}}{[\varpi_1,\lambda\delta]}+
 \underset{\substack{\phantom{\Vert}\\\deleted{\unsure{-}}\,2\varpi_0(\lambda\delta)}}{[\varpi_0,\lambda\delta]}\bigr] \]

\[ \begin{aligned}
 c_{\varpi_*}(\underset{\substack{\Vert\\\lambda\delta}}{\omega_1})&=\overbrace{\delta\,d\lambda}^{2}+\overbrace{\underset{\substack{\Vert\\\lambda\delta}}{\omega_1}(\overbrace{\varpi_1}^{2}-2\varpi_0)}^{1}\\
 c_{\varpi_*}(\underset{\substack{\Vert\\\mu.1}}{\omega_0})&=\overbrace{d\mu}^{1}+\overbrace{\omega_0(\overbrace{-\varpi_2}^{2}+\varpi_0)}^{0}\\
 c_{\varpi_*}(\underset{\substack{\Vert\\\nu\delta^{-1}}}{\omega_{-1}})&=\overbrace{\delta^{-1}d\nu}^{0}+\overbrace{\omega_{-1}(\overbrace{2\varpi_2}^{1}-\varpi_1)}^{0}
 \end{aligned} \]

\[ \begin{aligned}
 \unsure{\mathfrak g|\mathcal G}&:\quad \Omega,\ \underset{\substack{\Vert\\\underline O_X}}{\Omega^0},\ \Omega^{-1}\\
 \mathfrak g\otimes\omega&:\quad\Omega\otimes\underline\omega\quad\underline\omega\quad\Omega^{-1}\otimes\underline\omega
 \end{aligned} \]

On trouve \unsure{les conclusions suivantes|la conclusion suivante|la condition suivante} (\unsure{indépendantes|indépend$^{\mathrm{tes}}$|indépendante} du choix des splittings \ldots !)

\uline{tte connexion $c$ satisfait} $c(\mathrm{Fil}_i\,\mathfrak g)\subset\mathrm{Fil}_{i+1}(\mathfrak g\otimes\underline\omega)$ par passage au quotient, il induit
\[ c_{i,i+1}:\ \deleted{\unsure{\mathrm{gr}(\mathfrak g)}}\,\mathrm{gr}_i\,\mathfrak g\longrightarrow\mathrm{gr}_{i+1}(\mathfrak g\otimes\omega)\simeq\mathrm{gr}_{i+1}(\mathfrak g)\otimes\omega \]
qui est \uline{linéaire}, de façon précise

\[ \begin{aligned}
 c_{1,2}:\mathrm{gr}_1\longrightarrow\mathrm{gr}_2,\quad \Omega\longrightarrow\omega,&\quad\text{i.e. }c_{12}\in\Gamma\,\Omega^{-1}\otimes\underline\omega\\
 c_{2,3}:\mathrm{gr}_2\longrightarrow\mathrm{gr}_3\quad \underline O_X\longrightarrow\Omega^{-1}\otimes\omega,&\quad\text{i.e. }c_{23}\in\Gamma\,\Omega^{-1}\otimes\underline\omega
 \end{aligned} \]

on a [$c_{12}=-2\varpi_0$, $c_{2,3}=\varpi_0$, donc ]
\[ \boxed{c_{1,2}=-2c_{2,3}} \]

\uline{NB} On trouve que la connexion respecte la filtration ssi $\underline{c_{2,3}}=0$ (d'où $c_{12}=0$) [cela signifie $\varpi_0=0$] -- \deleted{\unsure{celles qui la}}\inserted{ce sont aussi} celles qui respectent le Borel $\mathscr L\subset\mathfrak g$, ou encore \deleted{telles} qui respectent la section $\sigma$ de $P$, i.e. pour lesquelles $\sigma$ est horizontale. Elles \unsure{sont|se} \deleted{\unsure{déterminées}}\unsure{diffèrent de|différenciées} une section de $\mathscr L\otimes\omega$. \deleted{\unsure{On obtient localement $f$, on chois}}

\margin{Re \uline{NB} On peut aussi considérer $c^\Omega:\underset{\substack{\Vert\\\check{\mathscr L}\simeq\mathscr L\otimes\Omega^{-1}}}{\mathfrak g/\Omega}\longrightarrow\Omega^{-1}\otimes\underline\omega$ induit par $c$, induisant évidemment sur $\Omega^0\simeq\underline O_X\subset\mathscr L\otimes\Omega^{-1}$, sa connaissance est \uline{équivalente} à celle des $c_\Omega$ et 2 signes, ets dans le cas général (\unsure{peut-être en|précisément} car. 2 ? C'est pas \unsure{sur} plus précis en}

\uline{NB} Les connexions sur $\mathfrak g$ compatibles avec crochet \deleted{\unsure{(plus}} (si 2 inv. -- sinon il faut parler des connexions sur $P$, ou $G$ etc) forment un torseur sous $\mathfrak g\otimes\underline\omega$. Or
\[ \mathfrak g\otimes\underline\omega\longrightarrow \Omega^{-1}\otimes\underline\omega \]
d'où un torseur sous $\Omega^{-1}\otimes\underline\omega$. \uline{Ce dernier} est splitté canoniq$^{\mathrm{t}}$, i.e. on a
\[ \underline{\mathrm{Conn}}\xrightarrow{\varpi_0}\Omega^{-1}\otimes\underline\omega \]
compatible avec \unsure{ce|$\varpi_0$}, et si $c$ section de $\underline{\mathrm{Conn}}$, $\varpi_0(c)$ détermine les hom $\mathrm{gr}_i(\mathfrak g)\to\mathrm{gr}_{i+1}(\mathfrak g\otimes\underline\omega)$

\[ \begin{aligned}
 \Omega=\mathrm{gr}_1(\mathfrak g)&\xrightarrow[-2\varpi_0]{}\omega=\mathrm{gr}_2(\mathfrak g\otimes\Omega^{\unsure{1}})\\
 \Omega^0=\underline O_X=\mathrm{gr}_2(\mathfrak g)&\xrightarrow[\varpi_0]{}\Omega^{-1}\otimes\underline\omega=\mathrm{gr}_3(\mathfrak g\otimes\underline\omega)
 \end{aligned} \]

Si on considère les $c$ pour lesquels $\varpi_0(c)\in\Gamma\underset{\substack{\Vert\\\mathrm{Hom}(\Omega,\underline\omega)}}{\Omega^{-1}\otimes\underline\omega}$ est donné \inserted{$i$}, on donne une restriction du Gr. struct. de $\underline{\mathrm{Conn}}$ à $\mathscr L\otimes\underline\omega$, d'où torseur $\mathrm{Conn}_i$ sous $\mathscr L\otimes\underline\omega$. Or on a
\[ \mathscr L\otimes\underline\omega\longrightarrow\underline\omega \]
d'où un torseur associé sous $\underline\omega$ : les conn. sur $\mathfrak g$ compatibles avec $i$, \uline{modulo} \inserted{op. d'}ajouter une section de $\Omega\otimes\underline\omega\subset\mathscr L\otimes\underline\omega\subset\mathfrak g\otimes\underline\omega$. C'est aussi (\deleted{\unsure{raisonnant}}) le torseur des restrictions à $\Omega$ des conn. compatibles avec $i$, en ce cas des hom additifs op. diff. d'ordre \uline{1} (rel $S$

\[ \Omega\xrightarrow{c_\Omega}\mathscr L\otimes\omega \]
telles que la composée \deleted{$\unsure{\Omega\to\Omega\otimes\omega}$} \inserted{$\rightarrow\mathscr L\otimes\omega\longrightarrow\omega$} soit linéaire (et égale à $-2i$), et que l'on ait
\unsure{sc|i.e.|($\ast\ast$)}
\[ c_\Omega(fw)-fc_\Omega(w)=w\otimes df\qquad \begin{aligned}w&\in\Gamma\Omega\\f&\in\Gamma\underline O_X\end{aligned} \]

\margin{TSVP}

% archive sheet number: 89
